% Part: second-order-logic
% Chapter: sol-and-sets
% Section: cardinalities

\documentclass[../../../include/open-logic-section]{subfiles}

\begin{document}

\olfileid{sol}{set}{crd}

\olsection{د سټونو شمېرنيزې اندازې}

\begin{explain}
لکه څنګه چې د تعريف ساحې متناهي يا نامتناهي، !!{enumerable} يا
!!{nonenumerable} والے څرګندولے شو، همداسې د~$\Domain{M}$ د يوه
فرعي سټ متناهي يا نامتناهي، !!{enumerable} يا !!{nonenumerable}
والے هم تعريفولے شو۔
\end{explain}

\begin{prop}
فورمول $\fn{Inf}(X) \ident$
\begin{multline*}
\lexists[u][(\lforall[x][\lforall[y][(\eq[u(x)][u(y)] \lif
      \eq[x][y])]] \land {} \\\lforall[x][(X(x) \lif X(u(x)))] \land {} \\
\lexists[y][(X(y) \land \lforall[x][(X(x)
      \lif \eq/[y][u(x)])])])]
\end{multline*}
د متغيرونو د ارزښت‌ټاکنې~$s$ له مخې هله او يوازې هله صدق کوي
چې $s(X)$ نامتناهي وي۔
\end{prop}
\textbf{د ژباړې د نسخې سمون (OLSOL-007):} د سرچينې په فورمول کښې
دا شرط نۀ و چې تابع د ټاکلي سټ غړي بېرته هماغه سټ ته ولېږي۔
له دې شرطه پرته حتا يو متناهي سټ هم فورمول پوره کولے شي؛
دلته د سټ د تړلتيا شرط ورزيات شو۔

\begin{prop}
فورمول $\fn{Count}(X) \ident $
\begin{multline*}
\lnot \lexists[x][X(x)] \lor {} \\
\lexists[z][\lexists[u][(X(z) \land
    \lforall[x][(X(x) \lif X(u(x)))] \land {} \\ \lforall[Y][((Y(z) \land
      \lforall[x][(Y(x) \lif Y(u(x)))]) \lif X \subseteq Y)])]]
\end{multline*}
د متغيرونو د ارزښت‌ټاکنې~$s$ له مخې هله او يوازې هله صدق کوي
چې $s(X)$ !!{enumerable} وي۔
\end{prop}
\textbf{د ژباړې د نسخې سمون (OLSOL-008):} د سرچينې بڼه خالي سټ
نۀ رانغاړي، که څه هم خالي سټ د شمېر وړ دے۔ همداراز په هغې کښې
د هر هغه سټ څخه د عين يووالي غوښتنه شوې وه چې پيلنے غړے
لري او د تابعې پر وړاندې تړلے وي؛ ټوله تعريف ساحه ئې يو
متناقض مثال دے کله چې ټاکلے سټ يو رښتينے فرعي سټ وي۔
خالي حالت ورزيات او وروستۍ پايله د فرعي‌والي شرط شوه۔

د کانتور له قضيې پوهېږو چې !!{nonenumerable} سټونه شته، او د
نامتناهي اندازو بېلابېلې ډېرې مرتبې هم شته۔ د سټونو تيوري
د سټونو د اندازو بشپړ حساب جوړوي او سټونو ته نامتناهي
شمېرنيز عددونه ټاکي۔ طبيعي عددونه د متناهي سټونو اندازې
اندازه کوي۔ د !!{denumerable} سټونو اندازه لومړۍ نامتناهي
شمېرنيزه اندازه ده چې~$\aleph_0$ («الف‌صفر») بلل کېږي۔
ورپسې نامتناهي اندازه~$\aleph_1$ ده۔ دا هغه تر ټولو کوچنۍ
اندازه ده چې يو سټ ئې لرلے شي خو د شمېر وړ نۀ وي (يعنې
اندازه ئې~$\aleph_0$ نۀ وي)۔ «$X$ د~$\aleph_0$ اندازه لري»
داسې تعريفولے شو: $\fn{Aleph}_0(X) \liff \fn{Inf}(X) \land \fn{Count}(X)$۔
$X$ د~$\aleph_1$ اندازه هله لري چې پخپله د شمېر وړ نۀ وي،
خو هر هغه فرعي سټ ئې چې له ټاکلي سټ څخه په شمېرنيزه اندازه
کوچنے وي، متناهي يا د~$\aleph_0$ اندازه ولري۔ نو دا په
دې توګه د~$\aleph_0$ تر کچې د کوچنيو اندازو شرط دے۔ دا په
لاندې فورمول څرګندېږي:
$\fn{Aleph_1}(X) \ident \lnot \fn{Count}(X) \land
\lforall[Y][((Y \subseteq X \land \lnot \cardeq{X}{Y}) \lif
  (\lnot\fn{Inf}(Y) \lor \fn{Aleph}_0(Y)))]$۔ د~$\aleph_2$
اندازه هم په ورته ډول تعريفېږي۔
\textbf{د ژباړې د نسخې سمون (OLSOL-009):} د سرچينې شرط
پر ټول ټاکلي سټ هم لګېده او د لومړۍ ناشمېرېدونکې اندازې سټ ئې
په تېروتنې سره رداوه؛ بل خوا متناهي سټونه ئې هم منل۔
دلته يوازې تر ټاکلي سټ په شمېرنيزه اندازه کوچني فرعي
سټونه محدودېږي، او پخپله ټاکلے سټ بايد د شمېر وړ نۀ وي۔

يوه اندازه ځانګړې پاملرنه غواړي: د پيوستار شمېرنيزه اندازه۔
دا د~$\Pow{\Nat}$ اندازه ده، يا په معادل ډول د~$\Real$ اندازه۔
دا چې يو سټ د پيوستار اندازه لري، په دويمې درجې منطق کښې هم
څرګندېږي، خو لږ زيات کار غواړي۔

\end{document}
